%% ============================================================
%% 分片 33：原文第 11 节「$A$-Branes and ${\cal D}$-Modules」（原文 glshort.tex 起于第 9938 行）
%% 由 tools/harvmac2tex.py 从 arXiv 源码（hep-th/0604151v3）自动生成。
%% 只译英文散文；公式、编号、\label、\cite 一律不改。
%% ============================================================

\subsection{${\cal D}$-modules Corresponding to $A$-Branes}
\label{sec:dmab}

 

The existence of the c.c. brane means that every $A$-brane in
complex structure $K$ is a module for the sheaf of differential
operators on $K_\M^{1/2}$.  In fact, in general, if ${\cal B}$ and
${\cal B}'$ are any two branes, then the $({\cal B},{\cal
B}')$-strings form a module for the algebra of $({\cal B},{\cal
B})$ strings.  This idea is illustrated in \dunamp\ of section
第~\ref{sec:comptwo}~节, and is part of the usual axioms of open-closed
topological field theory \cite{ref142,ref143}. To the
extent that sheafification is possible, the $({\cal B},{\cal B})$
strings form a sheaf of algebras, not just an  algebra, and the
$({\cal B},{\cal B}')$ strings form a sheaf of modules for this
sheaf of algebras. This statement just means that the ring and
module structures can be defined for open strings that are regular
only in a suitable open set in the target space.

We want to apply this construction for the case that ${\cal
B}={\cal B}_{c.c.}$ is the canonical coisotropic brane.  In this
case, we claim that for any brane ${\cal B}$ that is of $A$-type
in complex structure $K$, there are no instanton corrections to
$({\cal B}_{c.c},{\cal B}')$ strings.  The argument also applies
to higher order topological couplings (such as cubic Yukawa
couplings) involving the c.c. brane. The absence of instantons can
be argued in much the same way as for strings ending entirely on
the c.c. brane. The relevant disk instantons have a part of their
boundary on the c.c. brane, and on this part of the boundary the
instanton must be constant. But then, by analyticity, the
instanton must be constant everywhere.

We also know, from our previous investigation,  what sort of
sheafification is possible.  We can associate an algebra and a
module to an open set in $\MH$ of the form $T^*{{\cal V}}$,
${{\cal V}}\subset \M$, but not necessarily to more general open
sets. So an $A$-brane on $\MH(G,C)$ in complex structure $K$ gives
a sheaf of modules for the sheaf of algebras ${\cal D}_{K_{
\EUM}^{1/2}}$ over $\M$.

Now we will give a few examples of this. (For some examples worked
out in detail of open string quantization involving coisotropic
branes, see \cite{ref144}.)

Our first example is an $A$-brane ${\cal B}'$ defined by the
condition $\varphi=0$, with trivial Chan-Paton bundle. This
submanifold is a copy of $\cM(G,C)$ and is Lagrangian in complex
structures $K$ and $J$ and complex in complex structure $I$. Thus
${\cal B}'$ is an example of a $(B,A,A)$-brane. (We used this
brane in section 第~\ref{sec:opprod}~节\ in analyzing the operator product
expansion of 't Hooft operators.) We take the flat Chan-Paton
connection on ${\cal B}'$ to be trivial, for simplicity. As
described in \eqref{eq:exactcc}, the c.c. brane has a Chan-Paton line bundle
${\cal N}$ that is topologically trivial, but endowed with a
non-trivial connection.

To compute the spectrum of open strings on the classical level,
one can reduce the supersymmetric sigma-model to supersymmetric
quantum mechanics, i.e. retain only the zero modes. In this
approximation, open-string states with $({\cal B}_{c.c.},{\cal
B}')$ boundary conditions are sections of the tensor product with
${\cal N}^{-1}$ of a vector bundle obtained by quantizing the
space of fermion zero modes. The connection on ${\cal N}$ is zero
when restricted to $\varphi=0$, so we can omit the factor of
${\cal N}^{-1}$. The space of fermion zero modes is the tangent
space to $\M\subset \MH$ (fermions associated with normal
directions to $\M$ obey opposite boundary conditions at the two
ends of a string and have no zero modes). When we quantize the
space of fermion zero modes, we get the spinors on $\M$. Viewing
$\M$ as a complex manifold, its spin bundle is the same as
$K_\M^{1/2}\otimes(\oplus_{i=0}^{{\rm dim}\,\M}\Omega^{0,j}(\M))$.
Here $\Omega^{0,j}$ is the sheaf of $(0,j)$-forms on $\M$.

The BRST operator or topological supersymmetry $Q$ is in this
situation simply the $\bar\partial$ operator acting on the
$(0,j)$-forms with values in $K_\M^{1/2}$.  The sheaf of physical
 $({\cal B}_{c.c.},{\cal B}')$ strings is therefore the sheaf of
 holomorphic
sections of\footnote{The $\Omega^j$-forms with $j>0$ do not contribute
to the cohomology over a small open set ${{\cal V}}$, so they can
be omitted here.} $K_\M^{1/2}$. This is of course a sheaf of
modules for the sheaf of rings ${\cal D}_{K_{ \EUM}^{1/2}}$;
indeed, it is the sheaf of modules by which this sheaf of rings
was defined. This construction gives, possibly, a more direct
explanation of why the sheaf of rings derived from the c.c. brane
is precisely ${\cal D}_{K_{ \EUM}^{1/2}}$, rather than ${\cal
D}_{\cal L}$ for some other ${\cal L}$.

It may appear that we have implicitly used again the time-reversal
symmetry $\EUT$ in claiming that quantization of the space of
fermion zero modes gives precisely the spinors on $\M$, rather
than spinors with values in some line bundle ${\cal L}$. Actually,
we can give an alternative argument for this point, though it is
an argument that uses another discrete symmetry of the theory. The
branes ${\cal B}_{c.c.}$ and ${\cal B}'$ are physically sensible,
unitary branes in the four-dimensional supersymmetric gauge theory
defined on $M=\Bbb{R}\times I\times C$; here $\Bbb{R}$ is the time
direction, and $I$ is an interval with boundary conditions at the
two ends defined respectively by ${\cal B}_{c.c.}$ and ${\cal
B}'$.    To define the theory on $M$, there is some twisting to
preserve supersymmetry in the compactification on $C$, but no
twisting that involves the time direction.  So along with the
topological supersymmetry $Q$, the theory on $M$ is also invariant
under the adjoint supersymmetry $Q^\dagger$ (they obey
$\{Q,Q^\dagger\}=H$, where $H$ is the Hamiltonian).  The physical
theory has a  ``charge conjugation'' symmetry that exchanges $Q$
and $Q^\dagger$.  Invariance under this symmetry implies that the
Chan-Paton bundle obtained by quantizing the fermion zero modes is
precisely the spin bundle of $\CM$, not the tensor product of the
spin bundle with an additional line bundle.\footnote{If $\pi_1(G)$ is
nontrivial, we can modify the construction by taking the
Chan-Paton bundle on the brane ${\cal B}'$ to be a flat line
bundle.  Then we get spinors on $\M$ with values in a flat line
bundle; as a special case of this, we get all the possible spin
structures on $\M$. Because of the relation of $D$-branes to
$K$-theory \cite{ref102}, the choice of spin structure of $\M$ is really
part of a careful description of the Chan-Paton bundle of the
brane ${\cal B}'$.} This type of argument was first made in
\cite{ref145}\ in quantizing solitons.

Note that in this computation it was important that the support of
${\cal B}'$ is not only a Lagrangian submanifold with respect to
$\omega_K$, but is also a complex submanifold with respect to the
complex structure $\hat N=I$ determined by the c.c. brane. Otherwise,
the topological supercharge or BRST operator would not reduce to the $\bar\partial$ operator
in complex structure $I$.

Our second example is the case that is important for the geometric
Langlands program: a brane of type ${\CMF}$, that is a brane
${\cal B}_{\CMF}$ supported on a fiber ${\CMF}$ of the Hitchin
fibration with a flat Chan-Paton line bundle.  This is a brane of
type $(A,B,A)$, so in particular it is an $A$-brane in complex
structure $K$. Therefore, it gives rise to a sheaf of modules for
${\cal D}_{K_{ \EUM}^{1/2}}$. We will explain at the end of
section 第~\ref{sec:genotics}~节, by a further elementary argument, how to
convert this to a sheaf of modules for ${\cal D}$, the sheaf of
ordinary differential operators on $\M$. So the brane of type
${\CMF}$ has the two key properties: it is a magnetic eigenbrane
because it is the $S$-dual of a zerobrane; and it gives rise to a
sheaf of ${\cal D}$-modules over $\M$.  These are the basic claims
of the geometric Langlands program.

It is difficult to explicitly describe the sheaf of ${\cal D}$ or
${\cal D}_{K_{ \EUM}^{1/2}}$-modules that comes from a brane of
type $\CMF$. But we can do this in the the abelian case $G=U(1)$.
In this special case, the sigma-model is a free field theory, and
the boundary conditions are linear, so the computation of the
spectrum and module structure of the open strings is
straightforward.

For $G=U(1)$, we can think of $A$ and $\phi$ as real one-forms
($A$ is only defined up to gauge transformations, of course, while
$\phi$ is gauge-invariant).  The Hitchin equations for $A$ and
$\phi$ decouple, and in complex structure $I$, the Hitchin moduli
space is a product of the Jacobian ${\rm Jac}(C)$ of $C$ (which is
the moduli space of topologically trivial holomorphic line bundles
on $C$) and the vector space ${\EUBB}=H^0(C,\Omega^1)$. ${\rm
Jac}(C)$ is a complex torus of dimension $g_C$, the genus of $C$,
and the Hitchin moduli space ${\rm Jac}(C)\times {\EUBB}$ can be
identified with its cotangent bundle. The Hitchin fibration is the
projection to ${\EUBB}$.

Given $p\in {\EUBB}$ and the corresponding fiber ${\CMF}_p\simeq
{\rm Jac}(C)$ of the Hitchin fibration, we can compute the space
of $({\cal B}_{c.c.},{\cal B}_{{\CMF}_p})$ strings much as in the
previous example. One difference is that although the Chan-Paton
line bundle $\cL\to{\CMF}_p$ of the brane ${\cal B}_{{\CMF}_p}$ is
still topologically trivial, we now allow an arbitrary unitary
flat connection on it. Another difference is that the restriction
of the Chan-Paton line bundle ${\cal N}$ of the c.c. brane to
${\CMF}_p$ is holomorphically nontrivial, in general. On the other
hand, the canonical bundle $K_{{\CMF}_p}$ is trivial as ${\CMF}_p$
is a torus. Consequently, BRST-invariant open string states with
ghost number zero are holomorphic sections of $T_p=\cL\otimes
{\cal N}^{-1}\vert_{{\CMF}_p}$.

\def\bz{{\bar z}}
\def\Im{{\rm Im}}
\def\bpartial{{\bar\partial}}
But how does the  sheaf of  $({\cal B}_{c.c.},{\cal
B}_{c.c.})$ strings acts on sections of $T_p$?  We will answer this question
next. We work on $\RR\times I$ with the left
boundary on the brane ${\cal B}_{c.c.}$ and the right boundary on
${\cal B}_{{\CMF}_p}$.  The fiber $p$ of the Hitchin fibration
is defined by $\phi=\upsilon$, where $\upsilon=\upsilon_z\,dz+\upsilon_\bz\,d\bz$ is a real
harmonic one-form. The Chan-Paton
line bundle ${\cal L}\to{\CMF}_p$ of the brane ${\cal
B}_{{\CMF}_p}$ has a flat unitary connection that can be
conveniently represented as % 式 (11.17)
\setcounter{equation}{16}
\begin{equation}\label{eq:zoncob}
\beta={1\over
{2\pi}}\int_C |d^2 z|\left(a_z \delta A_\bz+a_\bz \delta
A_z\right)
\end{equation}
 where $a=a_z\,dz+a_\bz\,d\bz$ is a real harmonic
one-form on $C$.   Similarly, the connection on the line bundle
${\cal N}$ on the c.c. brane is\footnote{To get this formula and the next one,
we use eqn. \eqref{eq:tomog}\ for $\omega_J$ and $\omega_K$.  In eqn. \eqref{eq:tomog}, ${\rm Tr}$ represents
a trace in the $N$-dimensional representation of $U(N)$, and $A$ and $\phi$ are understood
to be skew-hermitian.  Since here we consider $A$ and $\phi$ as real one-forms,
we must include a minus sign.} % 式 (11.18)
\setcounter{equation}{17}
\begin{equation}\label{eq:homifly}
\alpha=-{\Im\,\tau\over
2\pi}\int_C |d^2 z|\left(\phi_z \delta A_\bz+\phi_\bz \delta
A_z\right).
\end{equation}
 Finally, the symplectic form $\omega=({\Im}\,\tau)\omega_K$ is the
exterior derivative of a one-form % 式 (11.19)
\setcounter{equation}{18}
\begin{equation}\label{eq:nomigly}
\zeta=-{i\,
\Im\,\tau\over 2\pi}\int_C |d^2 z|\left(\phi_\bz \delta A_z-\phi_z
\delta A_\bz\right).
\end{equation}

The action of the $A$-model, up to $Q$-exact
terms, is an integral over the boundary of $\Sigma$. The contribution to the action from
the right
boundary, which we call $\partial_R\Sigma$, is  % 式 (11.20)
\setcounter{equation}{19}
\begin{equation}\label{eq:omitro}
\int_{\partial_R\Sigma}
\Phi^*(\zeta -i\beta)={i\,\Im\,\tau\over 2\pi}\int_{\partial_R\Sigma}
\left(\left(\upsilon_z-(\Im\,\tau)^{-1}a_z\right) \dot{A}_\bz
-\left(\upsilon_\bz+(\Im\,\tau)^{-1}a_\bz\right) \dot{A}_z\right)ds.
\end{equation}

Here we used that on the right boundary, $\phi=\upsilon$.
Also, $s$ is a ``time'' coordinate on $\Sigma=\Bbb{R}\times I$, and
$\dot A=\partial A/\partial s$.
 The left boundary $\partial_L\Sigma$, having the opposite orientation,
contributes
% 式 (11.21)
\setcounter{equation}{20}
\begin{equation}\label{eq:bofloc}
-\int_{\partial_L\Sigma} \Phi^*(\zeta-i\alpha)=-{i\,\Im\,\tau\over \pi}\int_{\partial_L\Sigma}
 \phi_z
\dot{A}_\bz\,ds.
\end{equation}

The supersymmetric string states come from zero modes along the string, so to determine
them, the distinction
between $\partial_L\Sigma$ and $\partial_R\Sigma$ is unimportant and we can just add
the two contributions to the action.
The total action is accordingly
% 式 (11.22)
\setcounter{equation}{21}
\begin{equation}\label{eq:nofloc}
{i\,\Im\,\tau\over \pi}\int \left(-\phi_z+{1\over
2}(\upsilon_z-(\Im\,\tau)^{-1}a_z)\right)\dot{A}_\bz \,ds-{i\,\Im\,\tau\over
2\pi}\int(\upsilon_\bz+(\Im\,\tau)^{-1}a_\bz) \dot{A}_z\,ds.
\end{equation}

We see that the action depends on the parameters
$\upsilon_z,\upsilon_\bz,a_z,a_\bz$ in the combinations
% 式 (11.23)
\setcounter{equation}{22}
\begin{equation}\label{eq:numiko}
\upsilon_z-(\Im\,\tau)^{-1}a_z,\ \upsilon_\bz+(\Im\,\tau)^{-1}a_\bz.
\end{equation}

The first term in \eqref{eq:nofloc}, upon quantization, tells us how $\phi_z$ acts:
% 式 (11.24)
\setcounter{equation}{23}
\begin{equation}\label{eq:xococox}
\phi_z\to -i\pi(\Im\,\tau)^{-1}\left({\delta\over\delta
A_\bz}+{i\,\Im\,\tau\over 2\pi}\upsilon_z-{i\over 2\pi} a_z\right).
\end{equation}

This is a covariant $\partial$ operator on a topologically trivial line bundle
over ${\rm Jac}(C)$. The second term indicates that the
wavefunctions are sections of a topologically trivial but holomorphically nontrivial
 line bundle
 $T_p\to{\rm Jac}(C)$, which is, in fact, the same line bundle $T_p=\cL\otimes
{\cal N}^{-1}\vert_{{\CMF}_p}$ that we identified
 before. Indeed, this term is a total derivative, so if
we absorb it into the initial and final state wavefunctions, we
find that the wavefunctions are no longer independent of $A_z$ but
are annihilated by the operator
% 式 (11.25)
\setcounter{equation}{24}
\begin{equation}\label{eq:bilbox}
{\delta\over\delta A_z}-{i\,\Im\,\tau\over 2\pi}\upsilon_\bz-{i\over
2\pi} a_\bz .
\end{equation}

This is a covariant $\bpartial$ operator on a trivial line bundle
over ${\rm Jac}(C)$. Obviously, these two operators define a flat
connection on ${\rm Jac}(C)$. It is unitary if and only if
$\upsilon=0$.

So in short, quantization of the
$({\cal B}_{c.c.},{\cal B}_{{\CMF}_p})$ strings
has given us a
${\cal D}$-module associated to a choice of complex flat connection
on a trivial line bundle over ${\rm Jac}(C)$.  On the other hand,
$S$-duality identifies this family of $A$-branes with the set of
all zerobranes on $\MH(U(1),C)$, which in complex structure $J$ is
the moduli space of flat $\Bbb{C}^*$-bundles over $C$.  In other
words, in the abelian case $S$-duality establishes a natural
correspondence between gauge-equivalence classes of flat $\CC^*$
connections on $C$ and those on ${\rm Jac}(C)$.

This correspondence can be seen directly.  A flat $\CC^*$
connection on a manifold $M$ is the same as a homomorphism from
$H_1(M,\Bbb{Z})$ to $\CC^*$. Thus it is sufficient to show that
$H_1(C,\Bbb{Z})$ is isomorphic to $H_1({\rm Jac}(C),\Bbb{Z})$. But
since ${\rm Jac}(C)\simeq H^1(C,\RR)/H^1(C,\Bbb{Z})$, this is the
same as proving that $H_1(C,\Bbb{Z})\simeq H^1(C,\Bbb{Z})$, which
is Poincar\'e duality.
